\chapter[Calibration]{Calibration and Imaging}
\label{Chap:Calibration}
This would be how to calibrate and the calibration results for the datasets, plus any special considerations we took.

\section{Our observations}
\begin{itemize}
    \item We are using VLA data in X-band, what does this mean and what are the specs for the VLA.
    \item These observations are of list minihalos and this setup, observation proposal.
    \item Original source of these minihalos is the G14 paper (among others).
    \item Table of minihalos + observing setup.
\end{itemize}

\section{How to calibrate a dataset}
\begin{itemize}
    \item We have the VLA calibration pipeline to remove stress from our lives.
    \item We have a non-neglible amount of RFI to add stress to our lives. 
    \item A variety of steps but first pain, what is RFI and what does it cause.
\end{itemize}

\section{Flagging data}
\begin{itemize}
    \item Initial testing showed that the pipeline flagging steps didn't work well for our data.
    \item AOFLAGGER worked much better and has a mostly out-of-box VLA script.
    \item How does AOFLAGGER work
    \item Integrating into a pipeline step to still use the pipeline
\end{itemize}

\section{Calibration pipeline}
\begin{itemize}
    \item Work through each calibration step explaining what it does and why it is important. 
    \item Final preparation for imaging is averaging down to further reduce noise.
    \item Notes about general trends from the calibration, with RFI in 10-12 GHz D config being bad, much of it being flagged.
\end{itemize}

\section{Special mentions}
\begin{itemize}
    \item Either here or during following imaging section talk about the 2-3 datasets we processed slightly differently due to bad fits.
    \item Describe the fitting process for the custom bandpass.
\end{itemize}

\section{Imaging}
\begin{itemize}
    \item When we make a raw image from the data, we still have a lot of artefacts from the point spread function, which comes from the beam.
    \item This means we need a more advanced imaging process to remove this to get a clean image. Implemented in a few different ways, we choose tclean in CASA.
    \item Tclean works by modelling the PSF, creating an empty model, find brightest point in residuals, model that as real, subtract model from residuals, repeat 20K times.
    \item This works mostly well, but we can improve results using self calibration.
\end{itemize}

\subsection{Self calibration}
\begin{itemize}
    \item The phase calibration is the most iffy part for the main calibration, so we can make a test image and check the model seems reasonable.
    \item Then use that model to make small phase corrections to our dataset. Then repeat that a few times and these corrections should converge since it is circular.
    \item This is the final calibration steps that we do, phase corrections 20 degrees $\rightarrow$ 1 degree after a few rounds. 
\end{itemize}

\subsection{Final cleaning}
\begin{itemize}
    \item Between self calibrations this we also check the datasets and do some final manual cleaning.
    \item Clipping the datasets to a sensible amp range, removing spikes, SPWs where there is small amounts of data left.
    \item Overview of the \texttt{assess\_spw} function used for this.
\end{itemize}
